\section{Independent review of the no-push command interface}
\label{sec:nopush-independent-review}

This review concerns the new command semantics, independently of the
literal-tape verification. Let $I(a,b)=4p(a,b)$, where $p$ is Cantor
pairing, and write $[a;w]=p(a,w)$. Its total projections satisfy
$\operatorname{hd}[a;w]=a$, $\operatorname{tl}[a;w]=w$, and
$\operatorname{hd}(0)=\operatorname{tl}(0)=0$.
The identity $[0;0]=0$ is therefore retained throughout; no free-list
constructor assumption is made.

The state is $\langle t,w,d,h\rangle$. The no-push controller retains
the original Shift:
\[
 S\langle t,w,d,h\rangle=
 \begin{cases}
  \mathsf{halt},&t+h=1,\\
  \langle h,\operatorname{tl}(w),d,\operatorname{hd}(w)\rangle,&t+h\ne1.
 \end{cases}
\]
For $y=d+x$, it uses
\[
 \mathrm{MP}_x\langle t,w,d,h\rangle=
 \begin{cases}
  \langle y,w,y,0\rangle,&I(t,y)=h,\\
  \langle0,w,y,0\rangle,&I(t,y)\ne h.
 \end{cases}
\]
Thus successful MP retains its consequent in T and D and preserves the
underlying stack. This differs both from the original controller, which
pushes a dummy zero on success, and from the cached controller, which
returns the consequent in H and clears D.

The unchanged AX interface is
\[
 \mathrm{AX}_{c,x,y}\langle t,w,d,h\rangle
 =\langle K(h+y,d+x,c),[q_1,\ldots,q_9,t;w],0,0\rangle.
\]
Here $q_1$ is the separately emitted diagonal head, $q_2,\ldots,q_7$
are the six logical heads in reverse emission order, $q_8$ is the
restricted implication axiom, and $q_9=U(d,h)$.
The gate and $q_1,\ldots,q_8$ are multiples of four; $q_9$ has residue
two and is at least two. The gate is zero when its guard fails.

\subsection{Two useful total operations}

On any store with $d=0$, both branches of $\mathrm{MP}_0$ have exactly
the same output:
\[
 \mathrm{MP}_0\langle t,w,0,h\rangle=\langle0,w,0,0\rangle.
\]
This holds for arbitrary $t,h$, without deciding whether the equality
test succeeds. In particular, a single command gives a uniform
\emph{Discard} operation on normalized stores. The ability to coalesce
these branches is a real change in the command algebra: the old
push-on-success controller does not have this contract.

Let $A_0=\mathrm{AX}_{0,0,0}$ and compose commands from left to right.
Then
\[
 \mathrm{Push}=A_0;S^9;\mathrm{MP}_0,
 \qquad
 \mathrm{Pop}=A_0;S^{11}
\]
have the exact contracts
\begin{align*}
 \mathrm{Push}\langle t,w,d,h\rangle
   &=\langle0,[t;w],0,0\rangle,\\
 \mathrm{Pop}\langle t,[a;w],d,h\rangle
   &=\langle t,w,0,a\rangle.
\end{align*}
Both are total, including $t=1$. To verify this, AX first clears its
visible cache and offset. The first nine Shifts only test the gate or
adjacent generated heads; the stated residues exclude a sum of one.
After those nine Shifts, the top is $q_8$, the cache is $q_9$, and the
stack is $[t;w]$, so the final $\mathrm{MP}_0$ gives Push.
For Pop, the tenth Shift leaves $q_9$ in T and the saved $t$ in H.
The eleventh tests $q_9+t\ge2$, and therefore continues even when the
saved top is the contradiction code. This last inequality is essential:
one must not substitute the cached-controller Pop proof, whose final
guard tests only the cache and can halt at $t=1$.

These arguments use projection identities, so they also cover $w=0$,
$t=0$, and arbitrarily many written zero entries that collapse under
Cantor pairing. They do not require a syntactically distinguished stack
bottom.

\subsection{Axioms, generalization, and modus ponens}

The offset operation can be taken to be $\mathrm{MP}_i$ for every
$i\ge0$ on $\langle0,w,0,0\rangle$:
\[
 \mathrm{MP}_i\langle0,w,0,0\rangle=\langle0,w,i,0\rangle.
\]
For $i=0$ both branches coincide; for $i>0$, $I(0,i)>0$ forces the
failure branch. Hence
\[
 \mathrm{Gen}_i=
 \mathrm{Push};\mathrm{MP}_i;S;\mathrm{AX}_{i,0,0};S^{10}
\]
maps $\langle a,w,0,0\rangle$ to
$\langle U(i,a),w,0,0\rangle$ for every $a$, including $a=1$.
The Shift before AX tests two zero registers and retrieves $a$;
AX then clears that cache before further testing.

Starting with Discard, the ordinary separate-head selectors remain
valid. For $1\le j\le7$, use
\[
 \mathrm{Discard};\mathrm{AX}_{A,D,H};S^{j+1};\mathrm{Pop}^{9-j}.
\]
It returns $q_j$ in T and leaves $w$ unchanged, with D and H zero.
The last prescribed Pop consumes the deliberately saved zero, not an
element of the arbitrary underlying stack. The restricted implication
axiom with $u\le a$ uses
\[
 \mathrm{Discard};\mathrm{MP}_u;
 \mathrm{AX}_{b,a-u,c};S^9;\mathrm{Pop},
\]
and the guarded gate uses
\[
 \mathrm{Discard};\mathrm{AX}_{c,b,a};\mathrm{Pop}^{10}.
\]
For a genuine gate axiom its guard is satisfied. The formulas generated
by the remaining heads are never code one, so these selectors do not
introduce an unaccounted early halt.

Suppose the recursive compiler already handles proofs of $a$ and
$I(a,b)$, uniformly over the incoming top and stack and with D and H
zero. Write their command words as $C_a$ and $C_{I(a,b)}$.
Compile the MP node by
\[
 \mathrm{Discard};\mathrm{Push};
 C_{I(a,b)};\mathrm{Push};C_a;\mathrm{Push};
 S;S;\mathrm{MP}_b;\mathrm{Pop}.
\]
Unless a premise computation has already halted, the three pushes
prepare $\langle0,[a,I(a,b),0;w],0,0\rangle$.
The next two Shifts halt if $a=1$, which is an acceptable early
refutation. Otherwise they reach
$\langle a,[0;w],0,I(a,b)\rangle$.
Successful MP then gives
$\langle b,[0;w],b,0\rangle$.
The total Pop contract finally yields
$\langle b,w,0,0\rangle$, including when $b=1$.
The dummy zero is necessary to clear the outgoing H while preserving
the arbitrary tail $w$; deleting it would invalidate the compiler.

Induction on a numerical proof therefore supplies a primitive-recursive
word $C_p$ that halts or returns its conclusion in T with the original
underlying stack and zero D,H. If $p$ is a refutation, then
\[
             A_0;C_p;S
\]
halts from \emph{every} incoming store. The first AX normalizes D,H;
the recursive compiler accepts its arbitrary resulting top and stack;
the last Shift sees $t+h=1$ if an earlier halt did not occur. This
uniformity is what permits the existing exhaustive scheduler to carry
its store between trials.

\subsection{The converse invariant and the PA boundary}

For soundness, annotate T, H, and the canonical finite live stack by
numerical proofs. The zero tail carries the fixed truth proof. AX
preserves the annotation using its axiom suppliers and generalization.
A successful no-push MP obtains the proof of $y$ from the annotated
$t$ and $I(t,y)=h$ and places it in T; retaining $y$ in D is harmless,
because D has no theoremhood requirement. A failed MP returns truth.
At a Shift halt, natural-number arithmetic gives either $(t,h)=(1,0)$
or $(t,h)=(0,1)$, so an annotated refutation is available.

Every displayed command identity and compiler induction concerns finite
syntax and natural arithmetic and is formalizable in PA. Together with
the inherited numerical-calculus calibration, exhaustive enumeration,
and the separately checked positive physical simulations, these
identities provide the two directions of the requested ordinary
mathematical PA theorem. This review does not assert an independent
end-to-end Lean/Rocq elaboration or close the repository's iPRA endpoint.
